\section{Related Work}
The migration of Intrusion Detection Systems (IDS) from centralized cloud architectures to the network edge represents a paradigm shift necessitated by the rapid expansion of IoT ecosystems. This section reviews recent, pivotal literature across three intersecting domains critical to modern security: Machine Learning at the Edge, Real-Time Data Streaming Architectures, and Explainable AI (XAI) in cybersecurity, identifying the persistent gap our framework addresses.

\subsection{Machine Learning for IoT and Edge Security}
Recent literature emphasizes deploying ML models closer to data sources to mitigate latency and bandwidth saturation. Early approaches focused on inherently lightweight algorithms, such as k-Nearest Neighbors (k-NN) and Naive Bayes, to comply with the strict resource constraints of Edge nodes [1]. However, these legacy algorithms consistently fail to recognize the complex, high-dimensional attack patterns present in modern datasets like CICIDS2017. To achieve higher accuracy, recent studies have attempted to push Deep Learning (DL) models to the edge. For instance, Kim et al. [2] and Tan et al. [3] successfully deployed Convolutional Networks for IoT traffic analysis; however, these implementations rely heavily on advanced, high-cost Edge servers (e.g., NVIDIA Jetson) equipped with dedicated GPUs. Deploying state-of-the-art anomaly detection on pervasive, low-power, and economical hardware like the Raspberry Pi remains a profoundly unresolved challenge due to extreme memory constraints, often resulting in systemic packet loss during high-volume inference.

\subsection{Real-Time Streaming Architectures in IDS}
The efficacy of an edge IDS is entirely contingent upon its ability to ingest and process network telemetry in real-time. Batch-processing paradigms are obsolete against zero-day and automated DDoS attacks. Recent architectural studies have increasingly adopted distributed stream-processing frameworks. Apache Kafka has emerged as the industry standard for high-throughput, low-latency telemetry ingestion [4], [5]. While Kafka excels at managing massive data pipelines in cloud-based Security Information and Event Management (SIEM) systems, its integration directly onto localized, resource-constrained Edge nodes (such as IoT gateways) is sparsely documented. Current literature frequently overlooks the delicate tuning required to prevent the stream processing engine itself from consuming the Edge device's finite RAM, thereby starving the ML inference model.

\subsection{Explainable AI (XAI) in Network Traffic Analysis}
As the cybersecurity domain begrudgingly accepts "black-box" models to combat sophisticated attacks, Explainable AI (XAI) has transitioned from a theoretical concept to an operational imperative. Security Operations Center (SOC) personnel require actionable intelligence—an understanding of \textit{why} an alert triggered—to execute precise mitigation [6]. Post-hoc interpretation methods, notably Local Interpretable Model-agnostic Explanations (LIME) and SHapley Additive exPlanations (SHAP), have been successfully applied to cloud-based IDS retrospectively [7], [8]. These studies demonstrate SHAP's mathematical superiority in providing consistent, global, and local feature attributions.

\subsection{Addressing the Research Gap}
A critical review of the current literature reveals a disjointed research landscape. The majority of contemporary studies pursue one of two mutually exclusive paths: they either deploy highly accurate, interpretable XAI models exclusively on resource-rich \textit{Cloud} servers [7], or they deploy heavily compressed ML models to the \textit{Edge}, entirely abandoning interpretability to save computational resources [3].

Our proposed \textbf{X-IDS} framework explicitly bridges this exact gap. We position our research at the intersection of these three domains by demonstrating that XAI and Edge computing are not mutually exclusive. Instead of treating XAI as a post-deployment luxury, we mathematically leverage SHAP as an aggressive, pre-deployment dimensionality reduction tool. This XAI-driven pruning enables a robust ML architecture to operate flawlessly on a standard Raspberry Pi 4 under the sustained load of an Apache Kafka data stream, while simultaneously providing real-time, human-readable interpretability for every localized security alert.
